\section{Appendix: Supplied-propensity proofs and comparison algebra}

The supplied-propensity results combine inverse weighting, removal of the constant effect, and a quadratic comparison of independent blocks. The construction in \cref{obj:def:oracle-handle} applies over the continuous-carrier class of \cref{obj:def:oracle-broad-model}. Its upper bounds and the common-propensity lower witnesses give the matching separation guarantees in \cref{obj:thm:oracle-frontier-broad-class,obj:thm:oracle-frontier-resolved}. The comparison algebra then distinguishes the value of supplied assignment information on the original model from the cost of outcome tails at matched primitive smoothness.

\subsection{Inverse weighting and the public error bounds}

Fix an admissible tuple \(v=(p,\alpha,\beta,\gamma)\) and a law \(P\in\mathcal M_v^{\mathrm{e,br}}\), as defined in \cref{obj:def:oracle-broad-model}. Throughout this section, \(\lambda\) denotes Lebesgue probability measure on \([0,1]\). The conditional-mean identity in \cref{obj:thm:oracle-frontier-broad-class} identifies the mean effect through the untruncated inverse-propensity score. The treatment probability cancels its inverse weight separately in each arm, leaving the difference between the treated and baseline conditional means. This applies the inverse-probability weighting principle of \citet[pp.~665--666]{Horvitz1952} to the conditional mean contrast.

Outcome clipping makes this identity usable under the conditional moment envelope. The score \(R_i\), cutoff \(T\), and public bounds \(b,\Lambda\) are those of \cref{obj:def:oracle-handle}:
\[
b=20T^{1-p},
\qquad
\Lambda=\frac{80T^{2-p}}{s},
\qquad
s=\lfloor n/2\rfloor.
\]
The bias bound accounts for clipping in both treatment arms. The covariance bound combines the clipped second-moment envelope with the bounded inverse weights and averaging over an evaluation block. These calculations use uniform design, overlap, and the conditional moment restriction in
\cref{obj:ass:uniform,obj:ass:overlap,obj:ass:raw-moment}.

The statistical roles of the two bounds are distinct. Increasing the cutoff reduces clipping bias through \(T^{1-p}\), while the covariance allowance grows through \(T^{2-p}\). At \(p=2\), the latter factor equals one. Effect smoothness in \cref{obj:ass:effect-smoothness} controls histogram approximation of the mean contrast. Thus the broader-class attainment result depends on the moment order and effect smoothness, with continuous propensity and arm-mean representatives supplied by \cref{obj:def:oracle-broad-model}.

\subsection{Removing the unknown constant}

The histogram representation in \cref{obj:def:projection} makes the constant-effect direction explicit. Let
\[
u=J^{-1/2}(1,\ldots,1)^T,
\qquad
C_J=I_J-uu^T,
\]
where \(I_J\) is the \(J\times J\) identity matrix and \(C_J\) is the orthogonal projection onto the complement of the constant direction. In the notation of \cref{obj:def:oracle-handle}, \(u\) is a unit vector and
\[
C_JF_J(x)=F_J(x)-u.
\]
Consequently, the centered block vectors and their inner product have the equivalent representations
\[
V_b=\frac1s\sum_{i\in\mathcal I_b}C_JF_J(X_i)R_i,
\qquad b\in\{1,2\},
\]
and
\[
\langle V_1,V_2\rangle
=
\frac1{s^2}
\sum_{i\in\mathcal I_1}\sum_{k\in\mathcal I_2}
\bigl\{\Pi_J(X_i,X_k)-1\bigr\}R_iR_k.
\]
These are representations of the statistic already specified in \cref{obj:def:oracle-handle}.

Constant centering targets the distance of the effect from its average. Every unknown constant in the null of \cref{obj:def:oracle-broad-null} occupies the same removed direction. Clipping bias remains covered by the public bound \(b\), while effect smoothness governs the approximation of the remaining heterogeneity. The two disjoint evaluation blocks provide independent score vectors, so their quadratic comparison measures the squared population signal subject to the stated bias and covariance allowances.

The construction also specifies the test on its full supplied-function domain. For an arbitrary supplied continuous function \(f\), define its overlap-clipped version by
\[
\widetilde f(x)=\max\{1/4,\min\{f(x),3/4\}\}.
\]
The score \(R_i(f)\) and block vector \(V_b(f)\) denote the expressions in \cref{obj:def:oracle-handle} evaluated with \(\widetilde f\) in place of \(\widetilde e\). This convention bounds the denominators throughout the test domain. For a legal true supplied propensity, it agrees with the propensity itself, as required by the rejection expectation in \cref{obj:def:oracle-broad-risk}.

\subsection{Calibration, attainment, and common-propensity witnesses}

The threshold in \cref{obj:def:oracle-handle} is
\[
2b^2+1024\sqrt J\,\Lambda.
\]
Its first term accommodates the population discrepancy caused by clipping under a constant effect. The second supplies the quadratic fluctuation allowance. The resulting calibration in \cref{obj:thm:oracle-frontier-broad-class} controls rejection probability by \(0.01\) uniformly over the entire broader-class null for every \(n\ge2\). At the smallest sample sizes, the public tuning gives \(J=1\), for which constant centering makes the statistic identically zero.

The tuning coordinates clipping, projection, and quadratic variability. With the target scale \(a=s^{-E_0}\) from \cref{obj:def:oracle-handle}, the cutoff and rank are upper dyadic roundings of
\[
a^{-1/(p-1)}
\quad\text{and}\quad
a^{-1/\gamma},
\]
respectively. Their common separation exponent is
\[
E_0(p)=\frac{2\gamma q}{2\gamma+q},
\qquad q=\frac{p-1}{p},
\]
as defined in \cref{obj:def:sharp-frontier-scales}. The attained scale is \(\rho_n^{\mathrm e}(v)=n^{-E_0(p)}\). At moment order two,
\[
E_0(2)=\frac{2\gamma}{4\gamma+1}.
\]
This agrees numerically with the familiar one-dimensional quadratic-testing exponent in the Gaussian Sobolev setting of \citet[Section~6.1, pp.~249--250]{IngsterSapatinas2009GoodnessOfFit}, where the stated smoothness condition is strictly above \(1/4\). The supplied-propensity guarantees here have the conditional-moment and continuous-carrier scope specified in \cref{obj:thm:oracle-frontier-broad-class}, including \(\gamma=1/4\).

The lower and upper multipliers \(c_v^{\mathrm e},C_v^{\mathrm e}\) and threshold \(N_v^{\mathrm e}\) are defined in \cref{obj:thm:oracle-frontier-broad-class}. Their radius bounds hold for every \(n\ge2\). The power guarantee applies when
\[
C_v^{\mathrm e}\rho_n^{\mathrm e}(v)<D_v^{\mathrm{e,br}}.
\]
The public threshold places this upper separation at most \(d_0/2\), where \(d_0=1/(8\sqrt2)\) is the positive witness distance in that result. The lower bound \(D_v^{\mathrm{e,br}}\ge d_0\) then supplies a nonempty separation range.

The common-propensity witnesses connect this attainment result to both supplied experiments. In \cref{obj:thm:oracle-frontier-broad-class}, the null law and every positive-weight alternative law belong to the original and broader classes and share the propensity \(x\mapsto1/2\). Their complete-record distributions have total variation below one half. Supplying that same function preserves the comparison underlying the lower error bound.

Class inclusion transfers the broader-class calibration and power guarantees to the original model:
\[
\mathcal M_v\subseteq\mathcal M_v^{\mathrm{e,br}},
\qquad
H_0(v)\subseteq H_0^{\mathrm{e,br}}(v),
\]
using \cref{obj:def:model,obj:def:null,obj:def:oracle-broad-model,obj:def:oracle-broad-null}. The witnesses also lie within the original class, where \cref{obj:thm:oracle-frontier-resolved} supplies the corresponding necessity statement. Its multipliers \(c^{\mathrm e}_v,C^{\mathrm e}_v\) have the same displayed values as their broader-class counterparts; its power condition uses the original supremum \(D_v\).

\subsection{The preservation boundary on the original class}

Hold \(w=(\alpha,\beta,\gamma)\) fixed when evaluating the exponents in \cref{obj:def:sharp-frontier-scales}. The comparison identity in \cref{obj:thm:heavy-tail-comparison} is
\[
E_4(p)-E_0(p)
=
\frac{2\gamma q}{(2\gamma+q)D_p}\{F(p)-1\}.
\]
Its positive prefactor makes \(F(p)=1\) the equality boundary between the two separation exponents. The original-model exponent is their minimum.

The nonnegative comparison exponent \(G(v)\), defined in \cref{obj:thm:exact-propensity-preservation-boundary}, therefore has the representations
\[
G(v)=E_0(p)-E(p)
=
\frac{2\gamma q}{(2\gamma+q)D_p}
\max\{0,1-F(p)\}.
\]
The same-class radius ratio is bounded on both sides by positive multiples of \(n^{G(v)}\), with the exact multipliers given in that result. For \(F(p)\ge1\), this exponent equals zero and the ratio remains bounded over all \(n\ge2\). Equality points \(F(p)=1\) are included. For \(F(p)<1\), the positive exponent quantifies the polynomial gain from supplied assignment information. This yields precisely the preservation region in \cref{obj:def:preservation-region,obj:thm:exact-propensity-preservation-boundary}.

\subsection{The finite-moment and bounded comparison}

The bounded comparison evaluates the original-record exponent at moment order two while preserving \(w\). For presentation, let \(\delta_0(p)\) and \(\delta_4(p)\) denote the within-mechanism exponent differences. The formulas in \cref{obj:def:sharp-frontier-scales} give
\[
\delta_0(p)
:=E_0(2)-E_0(p)
=
\frac{4\gamma^2(2-p)}
{p(4\gamma+1)(2\gamma+q)},
\]
and
\[
\delta_4(p)
:=E_4(2)-E_4(p)
=
\frac{2S\beta(1/q-2)}{D_2D_p},
\qquad
D_2=1+2S+\frac{S}{2\gamma}.
\]
Both differences are strictly positive for \(p<2\) and vanish at \(p=2\).

At fixed smoothness, \(F(p)\ge F(2)\). The selected exponent difference consequently has three compatible representations:
\[
E(2)-E(p)=
\begin{cases}
\delta_0(p), & F(2)\ge1,\\[3pt]
\displaystyle
\delta_4(p)+
\frac{2\gamma q}{(2\gamma+q)D_p}\{F(p)-1\},
& F(p)\ge1>F(2),\\[8pt]
\delta_4(p), & F(p)<1.
\end{cases}
\]
At \(F(p)=1\), the additional term in the middle expression vanishes. At \(F(2)=1\), the bounded exponents coincide, so the adjoining expressions agree. These equality points preserve the exponent comparison across the change in the rate-determining mechanism.

The matched capped-radius inequalities in \cref{obj:thm:heavy-tail-comparison} attach the factor \(n^{E(2)-E(p)}\) to \(\Delta_n(v)\), the ratio defined in \cref{obj:def:comparison-handle}. Strict positivity for \(p<2\) gives the diverging ratio characterized by \cref{obj:def:tail-region}. At \(p=2\), the exponent difference is zero and the ratio remains uniformly bounded over sample sizes, as characterized by \cref{obj:def:equality-region}. The phase equality \(F(p)=1\) and the moment-order equality \(p=2\) thus describe distinct comparisons.

All these statements retain the finite-sample multipliers and saturation conventions of the cited radius results. Their local uniformity is expressed through common positive lower multipliers, finite upper multipliers, and finite bounds on the public power thresholds over compact subsets of the admissible parameter domains, as stipulated in \cref{obj:thm:heavy-tail-comparison,obj:thm:oracle-frontier-broad-class,obj:thm:oracle-frontier-resolved,obj:thm:exact-propensity-preservation-boundary}. These uniform constants accompany the parameter-dependent powers, including at the stated equality boundaries.

\begin{theoremv}{thm:heavy-tail-comparison}[Heavy-tail critical-radius comparison]
Let the admissible exponent domains be
\[
\begin{aligned}
\mathcal V
&=\{(p,\alpha,\beta,\gamma)\in\mathbb R^4:
1<p\le2,\ 0<\alpha\le1,\ 0<\beta\le1,\ 1/4\le\gamma\le1\},\\
\mathcal W
&=\{(\alpha,\beta,\gamma)\in\mathbb R^3:
0<\alpha\le1,\ 0<\beta\le1,\ 1/4\le\gamma\le1\}.
\end{aligned}
\]
For \(v=(p,\alpha,\beta,\gamma)\in\mathcal V\), write
\(w=(\alpha,\beta,\gamma)\). Use the public scales, multipliers, thresholds, and original-record tests defined in \cref{obj:def:sharp-frontier-scales}, and the models, null classes, risks, and capped radii defined in \cref{obj:def:model,obj:def:null,obj:def:bounded-model,obj:def:bounded-null,obj:def:testing-risk,obj:def:bounded-risk,obj:def:critical-radius,obj:def:bounded-radius,obj:def:binary-radius}. The matched radius ratio from \cref{obj:def:comparison-handle} is
\[
\Delta_n(v)=\frac{r_n^*(v)}{r_n^{*,\mathrm b}(w)}.
\]
Here \(\lambda\) is Lebesgue probability measure on \([0,1]\), and the population distance of the mean effect from its average is
\[
d(P)=
\left(
\int_0^1
\left|\tau_P(x)-\int_0^1\tau_P(z)\,d\lambda(z)\right|^2
\,d\lambda(x)
\right)^{1/2}.
\]
The maximal distances over the respective models are
\[
D_v=\sup_{P\in\mathcal M_v}d(P),
\qquad
D_w^{\mathrm b}=\sup_{P\in\mathcal M^{\mathrm b}_w}d(P),
\]
and \(d_0>0\) is the fixed witness distance supplied by \cref{obj:lem:causal-nonempty}. For the original-record experiment in \cref{obj:def:original-record-experiment}, write the rejection expectation as
\[
\mathsf R_n(P,\phi)=\mathbb E_{\mathsf S_{n,P}}[\phi].
\]

The following assertions hold.
\begin{itemize}
\item \textbf{(Matched radii and lower separations.)}
For every \(v\in\mathcal V\), \(c_v>0\) and \(C_v>0\). For every integer \(n\ge2\), both scales are positive and
\[
\begin{aligned}
c_v\rho_n(v)
&\le r_n^*(v)\le C_v\rho_n(v),\\
c_w^{\mathrm b}\rho_n^{\mathrm b}(w)
&\le r_n^{*,\mathrm b}(w)
=r_n^{*,\mathrm{bin}}(w)
\le C_w^{\mathrm b}\rho_n^{\mathrm b}(w).
\end{aligned}
\]
The lower separations satisfy
\[
0<c_v\rho_n(v)<D_v,
\qquad
0<c_w^{\mathrm b}\rho_n^{\mathrm b}(w)<D_w^{\mathrm b},
\]
and the corresponding minimax type-II errors satisfy
\[
B_n\!\left(v,c_v\rho_n(v)\right)>\frac1{10},
\qquad
B_n^{\mathrm b}\!\left(w,c_w^{\mathrm b}\rho_n^{\mathrm b}(w)\right)>\frac1{10}.
\]

\item \textbf{(Size and conditional power.)}
For every \(v\in\mathcal V\) and every integer \(n\ge2\),
\[
\begin{aligned}
\mathsf R_n(P,\phi^*_{n,v})&\le\frac1{10}
&&\text{for every }P\in H_0(v),\\
\mathsf R_n(P,\phi^{*,\mathrm b}_{n,w})&\le\frac1{10}
&&\text{for every }P\in H_0^{\mathrm b}(w).
\end{aligned}
\]
If \(C_v\rho_n(v)<D_v\), then
\[
1-\mathsf R_n(P,\phi^*_{n,v})\le\frac1{10}
\]
for every \(P\in\mathcal M_v\) satisfying \(d(P)\ge C_v\rho_n(v)\).
If \(C_w^{\mathrm b}\rho_n^{\mathrm b}(w)<D_w^{\mathrm b}\), then
\[
1-\mathsf R_n(P,\phi^{*,\mathrm b}_{n,w})\le\frac1{10}
\]
for every \(P\in\mathcal M^{\mathrm b}_w\) satisfying
\(d(P)\ge C_w^{\mathrm b}\rho_n^{\mathrm b}(w)\).

\item \textbf{(Vanishing scales and public thresholds.)}
For every \(v\in\mathcal V\),
\[
\rho_n(v)\longrightarrow0,
\qquad
\rho_n^{\mathrm b}(w)\longrightarrow0
\quad\text{as }n\longrightarrow\infty.
\]
For every natural number \(n\),
\[
\begin{aligned}
n\ge N_v
&\ \Longrightarrow\ C_v\rho_n(v)<d_0,\\
n\ge N_w^{\mathrm b}
&\ \Longrightarrow\ C_w^{\mathrm b}\rho_n^{\mathrm b}(w)<d_0.
\end{aligned}
\]

\item \textbf{(Phase identities and radius comparison.)}
Use the supplied-propensity exponent \(E_0(p)\), rough-confounding exponent \(E_4(p)\), phase-comparison scalar \(F(p)\), and denominator \(D_p\) from \cref{obj:def:sharp-frontier-scales}, with the smoothness tuple \(w\) fixed. With
\[
q=\frac{p-1}{p},
\qquad
E(p)=\min\{E_0(p),E_4(p)\},
\]
every \(v\in\mathcal V\) satisfies
\[
E_4(p)-E_0(p)
=
\frac{2\gamma q}{(2\gamma+q)D_p}\bigl(F(p)-1\bigr),
\]
and
\[
F(p)\ge1\ \Longrightarrow\ E(p)=E_0(p),
\qquad
F(p)<1\ \Longrightarrow\ E(p)=E_4(p).
\]
With
\[
D_2=1+2S+\frac{S}{2\gamma},
\qquad
F_2=2S+\frac{S}{2\gamma},
\]
the exponent difference, with evaluation at \(2\) preserving \(w\), is
\[
E(2)-E(p)=
\begin{cases}
\displaystyle
\frac{4\gamma^2(2-p)}{(4\gamma+1)(2\gamma p+p-1)},
&F_2\ge1,\\[6pt]
\displaystyle E_4(2)-E_0(p),
&F_2<1\le F(p),\\[4pt]
\displaystyle
\frac{2S\beta(2-p)}{(p-1)D_pD_2},
&F(p)<1.
\end{cases}
\]
For every integer \(n\ge2\),
\[
\frac{c_v}{C_w^{\mathrm b}}\,n^{E(2)-E(p)}
\le\Delta_n(v)
\le
\frac{C_v}{c_w^{\mathrm b}}\,n^{E(2)-E(p)},
\qquad
1\le\Delta_n(v).
\]

\item \textbf{(Strict and equality regions.)}
The strict-enlargement and bounded-ratio regions of \cref{obj:def:tail-region,obj:def:equality-region} are exactly
\[
\mathcal T=\{v\in\mathcal V:p<2\},
\qquad
\mathcal E=\{v\in\mathcal V:p=2\}.
\]
There exists a set \(U\subseteq\mathbb R^4\), open in the product topology of the four public exponents, such that
\[
U\cap\mathcal V\ne\varnothing,
\qquad
U\cap\mathcal V\subseteq\mathcal T.
\]
For every \(v\in\mathcal V\) with \(p<2\),
\[
E(p)<E(2),
\qquad
\frac{\rho_n(v)}{\rho_n^{\mathrm b}(w)}
\longrightarrow\infty
\quad\text{as }n\longrightarrow\infty.
\]

\item \textbf{(Tail-essential complete-record priors.)}
For every \(v\in\mathcal V\) with \(p<2\), the selected mark magnitudes from \cref{obj:def:converse-handle} satisfy
\[
L_{n,v}\longrightarrow\infty
\quad\text{as }n\longrightarrow\infty.
\]
For every integer \(n\ge2\), the selected priors
\(\pi_{0,n,v}\) and \(\pi_{1,n,v}\) from \cref{obj:def:converse-handle} are normalized finite probability priors. Every law receiving positive null-prior weight belongs to \(H_0(v)\), and every law receiving positive alternative-prior weight belongs to \(\mathcal M_v\) and satisfies
\[
c_v\rho_n(v)\le d(P)<D_v.
\]
For every law receiving positive weight in either prior, each conditional arm law assigns probability one to the three-point mark set: writing \(Q_{a,P}(dy\mid x)\) for the conditional outcome probability kernel in arm \(a\),
\[
Q_{a,P}\!\left(\{0,-L_{n,v},L_{n,v}\}\mid x\right)=1
\quad
\text{for }a\in\{0,1\}\text{ and }\lambda\text{-almost every }x.
\]
The complete-record mixtures defined in \cref{obj:def:converse-handle} satisfy
\[
\operatorname{TV}\!\left(\mathbb P_{0,n,v},\mathbb P_{1,n,v}\right)<\frac12,
\]
where total variation is the supremum of the absolute difference in event probabilities.

For each such \(v\), there exists a natural-number threshold \(N(v)\) such that, for every \(n\ge N(v)\), every pair of normalized finite probability priors supported respectively on
\[
H_0^{\mathrm b}(w)
\quad\text{and}\quad
\{P\in\mathcal M^{\mathrm b}_w:c_v\rho_n(v)\le d(P)\}
\]
has complete-record mixture total variation at least \(1/2\), in the sense of \cref{obj:thm:tail-essential-full-record-converse}.

\item \textbf{(Compact uniformity.)}
For every compact set of admissible public tuples, there exist real constants \(c,C,N\), with \(c>0\), such that
\[
c\le c_v,\qquad C_v\le C,\qquad N_v\le N
\]
for every tuple in that set. Likewise, for every compact set of admissible smoothness tuples, there exist real constants \(c,C,N\), with \(c>0\), such that
\[
c\le c_w^{\mathrm b},
\qquad
C_w^{\mathrm b}\le C,
\qquad
N_w^{\mathrm b}\le N
\]
for every smoothness tuple in that set.
\end{itemize}
\end{theoremv}

\begin{theoremv}{thm:oracle-frontier-broad-class}[Supplied-propensity broader-class frontier]
For every admissible public tuple \(v=(p,\alpha,\beta,\gamma)\) of \cref{obj:def:model}, consider the broader model and its constant-effect null in \cref{obj:def:oracle-broad-model,obj:def:oracle-broad-null}, with testing risk and capped critical radius as in \cref{obj:def:oracle-broad-risk,obj:def:oracle-broad-radius}. Put
\[
q=\frac{p-1}{p},\qquad
E_0=\frac{2\gamma q}{2\gamma+q},\qquad
\rho_n^{\mathrm e}(v)=n^{-E_0},\qquad
d_0=\frac{1}{8\sqrt2},
\]
and define the lower multiplier, upper multiplier, and public sample-size threshold by
\[
c_v^{\mathrm e}=\frac{4^{-\gamma}}{16\sqrt3},\qquad
C_v^{\mathrm e}=3\bigl(120+128\sqrt{160\sqrt2}\bigr),\qquad
N_v^{\mathrm e}
=\left\lceil\max\left\{2,
\left(\frac{2C_v^{\mathrm e}}{d_0}\right)^{1/E_0}
\right\}\right\rceil .
\]
Let \(\lambda\) denote Lebesgue probability measure on \([0,1]\), and write the heterogeneity distance as
\[
d(P)=
\left\{\int_{[0,1]}
\left(\tau_P(x)-\int_{[0,1]}\tau_P(z)\,\lambda(dz)\right)^2
\,\lambda(dx)\right\}^{1/2}.
\]
Then
\[
d_0\le D_v^{\mathrm{e,br}}\le\frac12,\qquad
\frac{1}{64\sqrt3}\le c_v^{\mathrm e},\qquad
0<C_v^{\mathrm e}.
\]

For every integer \(n\ge2\), the following guarantees hold.
\begin{itemize}
\item \textbf{(Critical radius.)}
\[
c_v^{\mathrm e}\rho_n^{\mathrm e}(v)
\le r_n^{*,\mathrm{e,br}}(v)
\le C_v^{\mathrm e}\rho_n^{\mathrm e}(v),
\qquad
B_n^{\mathrm{e,br}}
\bigl(v,c_v^{\mathrm e}\rho_n^{\mathrm e}(v)\bigr)
\ge\frac25.
\]

\item \textbf{(Calibration and power.)}
The supplied-propensity test \(\phi^{*,\mathrm e}_{n,v}\) of
\cref{obj:def:oracle-handle} satisfies
\[
\mathsf R_n^{\mathrm e}
(P,\phi^{*,\mathrm e}_{n,v})\le0.01
\qquad
\text{for every }P\in H_0^{\mathrm{e,br}}(v).
\]
Here \(\mathsf R_n^{\mathrm e}(P,\phi^{\mathrm e})\) is the expected rejection probability when the test receives the entire true continuous propensity \(e_P\), together with the independent original-record sample and public randomization. If
\[
C_v^{\mathrm e}\rho_n^{\mathrm e}(v)<D_v^{\mathrm{e,br}},
\]
then every \(P\in\mathcal M_v^{\mathrm{e,br}}\) with
\(d(P)\ge C_v^{\mathrm e}\rho_n^{\mathrm e}(v)\) satisfies
\[
1-\mathsf R_n^{\mathrm e}
(P,\phi^{*,\mathrm e}_{n,v})\le0.01.
\]

\item \textbf{(Finite lower-bound witnesses.)}
There exist a continuous function \(f\), a null law \(P_0\), and a normalized finite probability prior
\[
f(x)=\frac12\quad(x\in[0,1]),\qquad
\pi_1=\sum_{j=1}^{k}\omega_j\delta_{P_j},\qquad
\omega_j\ge0,\qquad
\sum_{j=1}^{k}\omega_j=1,
\]
with at least one strictly positive weight, such that
\[
P_0\in H_0(v)\cap H_0^{\mathrm{e,br}}(v),
\qquad e_{P_0}=f.
\]
Here \(H_0(v)\) is the original-model null of \cref{obj:def:null}.
For every \(j\) with \(\omega_j>0\),
\[
P_j\in\mathcal M_v\cap\mathcal M_v^{\mathrm{e,br}},
\qquad e_{P_j}=f,\qquad
c_v^{\mathrm e}\rho_n^{\mathrm e}(v)\le d(P_j)<d_0,
\]
where \(\mathcal M_v\) is defined in \cref{obj:def:model}.
Define the complete-record alternative mixture by
\[
\overline P_1^{(n)}
=\sum_{j=1}^{k}\omega_j P_j^{\otimes n}.
\]
With total variation defined as the supremum of absolute event-probability differences,
\[
\operatorname{TV}
\bigl(P_0^{\otimes n},\overline P_1^{(n)}\bigr)<\frac12.
\]
For every allowed randomized supplied-propensity test \(\phi^{\mathrm e}\) satisfying
\[
\mathsf R_n^{\mathrm e}(P,\phi^{\mathrm e})\le\frac1{10}
\qquad
\text{for every }P\in H_0^{\mathrm{e,br}}(v),
\]
there is an index \(j\) with \(\omega_j>0\) such that
\[
1-\mathsf R_n^{\mathrm e}(P_j,\phi^{\mathrm e})
\ge\frac25.
\]
\end{itemize}

For every integer \(n\ge N_v^{\mathrm e}\),
\[
C_v^{\mathrm e}\rho_n^{\mathrm e}(v)\le\frac{d_0}{2}.
\]

For each \(P\in\mathcal M_v^{\mathrm{e,br}}\), let
\(Q_{a,P}(dy\mid x)\) denote the conditional outcome probability kernel in treatment arm \(a\). The untruncated inverse-propensity score has conditional mean equal to the canonical effect: for \(\lambda\)-almost every \(x\),
\[
e_P(x)\int_{\mathbb R}\frac{y}{e_P(x)}\,Q_{1,P}(dy\mid x)
+
\bigl(1-e_P(x)\bigr)
\int_{\mathbb R}\frac{-y}{1-e_P(x)}\,Q_{0,P}(dy\mid x)
=\tau_P(x).
\]

The constants and exponent are locally uniform on compact subsets of the admissible public parameter domain, equipped with the product topology on \((p,\alpha,\beta,\gamma)\).
\end{theoremv}

\begin{theoremv}{thm:oracle-frontier-resolved}[Supplied-propensity frontier and attainment]
Let \(\mathcal V\) denote the admissible domain of public tuples \(v=(p,\alpha,\beta,\gamma)\) in \cref{obj:def:model}, equipped with the topology inherited from \(\mathbb R^4\). For \(v\in\mathcal V\), define
\[
q=\frac{p-1}{p},\qquad
E_0=\frac{2\gamma q}{2\gamma+q},\qquad
\rho_n^{\mathrm e}(v)=n^{-E_0},\qquad
c^{\mathrm e}_v=\frac{4^{-\gamma}}{16\sqrt3},\qquad
C^{\mathrm e}_v=3\left(120+128\sqrt{160\sqrt2}\right).
\]
Use the supplied-propensity rule \(\phi^{*,\mathrm e}_{n,v}\) of \cref{obj:def:oracle-handle}, and the risk and capped radius of \cref{obj:def:oracle-risk,obj:def:oracle-radius}. Write \(\mathsf R_n^{\mathrm e}(P,\phi^{\mathrm e})\) for the expected rejection probability when the entire true continuous propensity \(e_P\) is supplied to the test.

Let \(\lambda\) be Lebesgue probability measure on \([0,1]\), and define
\[
d(P)=
\left\{\int_0^1
\left(\tau_P(x)-\int_0^1\tau_P(z)\,d\lambda(z)\right)^2
\,d\lambda(x)\right\}^{1/2},
\qquad
D_v=\sup_{P\in\mathcal M_v}d(P).
\]
Let \(d_0>0\) be the fixed distance supplied by the legal model witness, and set
\[
N^{\mathrm e}_v=
\left\lceil
\max\left\{2,\left(\frac{2C^{\mathrm e}_v}{d_0}\right)^{1/E_0}\right\}
\right\rceil.
\]
Then
\[
\frac{1}{64\sqrt3}\le c^{\mathrm e}_v,
\qquad
0<C^{\mathrm e}_v.
\]
For every integer \(n\ge2\),
\[
c^{\mathrm e}_v\rho_n^{\mathrm e}(v)
\le r_n^{*,\mathrm e}(v)
\le C^{\mathrm e}_v\rho_n^{\mathrm e}(v),
\qquad
B_n^{\mathrm e}\!\left(v,c^{\mathrm e}_v\rho_n^{\mathrm e}(v)\right)
\ge\frac25.
\]
The attaining rule satisfies
\[
\mathsf R_n^{\mathrm e}\!\left(P,\phi^{*,\mathrm e}_{n,v}\right)
\le0.01
\qquad(P\in H_0(v)).
\]
Whenever \(C^{\mathrm e}_v\rho_n^{\mathrm e}(v)<D_v\), it also satisfies
\[
1-\mathsf R_n^{\mathrm e}\!\left(P,\phi^{*,\mathrm e}_{n,v}\right)
\le0.01
\qquad
\left(P\in\mathcal M_v,\quad
C^{\mathrm e}_v\rho_n^{\mathrm e}(v)\le d(P)\right).
\]

For each such \(v,n\), there exist a null law \(P_0\in H_0(v)\) and a normalized finite prior
\[
\pi_1=\sum_{j=1}^{k}\omega_j\delta_{P_j},
\qquad
\omega_j\ge0,\qquad
\sum_{j=1}^{k}\omega_j=1,\qquad
\exists j:\ \omega_j>0,
\]
such that every positive-weight support law satisfies
\[
P_j\in\mathcal M_v,\qquad
e_{P_j}=e_{P_0},\qquad
c^{\mathrm e}_v\rho_n^{\mathrm e}(v)\le d(P_j)<D_v.
\]
Define its complete-record mixture by
\[
\mathbb P_{\pi_1,n}=\sum_{j=1}^{k}\omega_jP_j^{\otimes n},
\]
and let \(\operatorname{TV}\) denote the supremum of absolute event-probability differences. These laws satisfy
\[
\operatorname{TV}\!\left(P_0^{\otimes n},\mathbb P_{\pi_1,n}\right)<\frac12.
\]
For every supplied-propensity test \(\phi^{\mathrm e}\) satisfying
\[
\mathsf R_n^{\mathrm e}(P,\phi^{\mathrm e})\le\frac1{10}
\qquad(P\in H_0(v)),
\]
there is an index \(j\) with
\[
\omega_j>0,\qquad
1-\mathsf R_n^{\mathrm e}(P_j,\phi^{\mathrm e})\ge\frac25.
\]
For every integer \(n\ge N^{\mathrm e}_v\),
\[
C^{\mathrm e}_v\rho_n^{\mathrm e}(v)\le\frac{d_0}{2}.
\]

For \(P\in\mathcal M_v\), let \(Q_{a,P}(dy\mid x)\) denote the conditional outcome probability kernel in treatment arm \(a\). The untruncated inverse-propensity score identifies the original mean effect: for \(\lambda\)-almost every \(x\),
\[
e_P(x)\int_{\mathbb R}\frac{y}{e_P(x)}\,Q_{1,P}(dy\mid x)
+
\bigl(1-e_P(x)\bigr)
\int_{\mathbb R}\frac{-y}{1-e_P(x)}\,Q_{0,P}(dy\mid x)
=\tau_P(x).
\]
The frontier constants and nonempty-power thresholds are locally uniform over admissible tuples: on every compact subset of \(\mathcal V\), the lower multipliers admit a common positive lower bound, the upper multipliers admit a common finite upper bound, and the thresholds admit a common finite upper bound, including at \(p=2\) and \(\gamma=1/4\).
\end{theoremv}

\begin{theoremv}{thm:exact-propensity-preservation-boundary}[Exact propensity preservation boundary]
Let \(\mathcal V\) be the domain of public tuples \(v=(p,\alpha,\beta,\gamma)\) satisfying:
\begin{itemize}
\item \textbf{(Moment order.)} \(1<p\le2\).
\item \textbf{(Propensity smoothness.)} \(0<\alpha\le1\).
\item \textbf{(Baseline smoothness.)} \(0<\beta\le1\).
\item \textbf{(Effect smoothness.)} \(1/4\le\gamma\le1\).
\end{itemize}
For every \(v\in\mathcal V\), use the quantities and original-model radius multipliers \(c_v,C_v\) of \cref{obj:def:sharp-frontier-scales}, with
\[
q=\frac{p-1}{p},\qquad S=\alpha+\beta,\qquad
D_p=1+2\alpha+\frac{\beta}{q}+\frac{S}{2\gamma},
\]
\[
E_0(p)=\frac{2\gamma q}{2\gamma+q},\qquad
E_4(p)=\frac{2S}{D_p},\qquad
F(p)=\frac{2\alpha}{p-1}+\frac{\beta}{q}+\frac{S}{2\gamma}.
\]
Define the comparison exponent and the supplied-propensity radius multipliers by
\[
G(v)=\max\{0,E_0(p)-E_4(p)\},\qquad
c^{\mathrm e}_v=\frac{4^{-\gamma}}{16\sqrt3},\qquad
C^{\mathrm e}_v=3\bigl(120+128\sqrt{160\sqrt2}\bigr).
\]
Then
\[
G(v)=
\frac{2\gamma q}{(2\gamma+q)D_p}\max\{0,1-F(p)\}.
\]
For the unknown- and supplied-propensity experiments on the same class \(\mathcal M_v\) of \cref{obj:def:model}, the capped critical radii of \cref{obj:def:critical-radius,obj:def:oracle-radius} satisfy, for every integer \(n\ge2\),
\[
\frac{c_v}{C^{\mathrm e}_v}n^{G(v)}
\le
\frac{r_n^*(v)}{r_n^{*,\mathrm e}(v)}
\le
\frac{C_v}{c^{\mathrm e}_v}n^{G(v)}.
\]
If \(F(p)=1\), then
\[
0<\frac{c_v}{C^{\mathrm e}_v},
\qquad
0<\frac{C_v}{c^{\mathrm e}_v},
\]
and, for every integer \(n\ge2\),
\[
\frac{c_v}{C^{\mathrm e}_v}
\le
\frac{r_n^*(v)}{r_n^{*,\mathrm e}(v)}
\le
\frac{C_v}{c^{\mathrm e}_v}.
\]
If \(F(p)<1\), then
\[
\lim_{n\to\infty}\frac{r_n^*(v)}{r_n^{*,\mathrm e}(v)}=+\infty,
\]
and, for every integer \(n\ge2\),
\[
\frac{c_v}{C^{\mathrm e}_v}n^{E_0(p)-E_4(p)}
\le
\frac{r_n^*(v)}{r_n^{*,\mathrm e}(v)}
\le
\frac{C_v}{c^{\mathrm e}_v}n^{E_0(p)-E_4(p)}.
\]
The preservation region of \cref{obj:def:preservation-region} is exactly
\[
\mathcal R=
\left\{v\in\mathcal V:
\frac{2\alpha}{p-1}+\frac{\beta}{q}
+\frac{\alpha+\beta}{2\gamma}\ge1\right\}.
\]
Moreover, for every compact subset \(K\) of \(\mathcal V\), in the product topology on the four real coordinates, there exist real constants \(c,C\), with \(c>0\), such that, for every \(v\in K\),
\[
c\le\frac{c_v}{C^{\mathrm e}_v},
\qquad
64\sqrt3\,C_v\le C.
\]
\end{theoremv}
